\documentclass{report}

\usepackage[dvipsnames]{xcolor}
\usepackage{listings}
\usepackage[utf8]{inputenc}
\usepackage{alloy-style}
\usepackage{textcomp}
\usepackage{amsmath}

% Package for pdf metadata
\usepackage{hyperref}
\usepackage{graphicx}
\usepackage{titlepic}

\usepackage[nottoc,notlof,notlot]{tocbibind} 
\renewcommand\bibname{References}

\hypersetup{
  pdftitle={RASD},
  pdfauthor={Carlo Dell'Acqua, Adriana Ferrari, Angelica Sofia Valeriani},                     % author
  pdfsubject={SafeStreet},
  pdfkeywords={RASD, SafeStreet, Software Engineering},
  colorlinks=true,
  linkcolor=blue,
  citecolor=black,
  filecolor=black,
  urlcolor=purple,
  linktoc=page
}

\setcounter{secnumdepth}{3}
  % Cover

\begin{titlepage}
        \titlepic{\includegraphics[width=9cm]{RASD_Images/logo.jpg}}
        \title{RASD}
        \author{Carlo Dell'Acqua, Adriana Ferrari, Angelica Sofia Valeriani}
\end{titlepage}

\begin{document}
\maketitle
  % Table of contents
  \tableofcontents

  % 1 - Introduction
  \chapter{Introduction}
  \definecolor{gray}{HTML}{F7F6ED}

\lstinputlisting[
  language=alloy,
  numbers=left,
  stepnumber=1,
  breaklines=true,
  backgroundcolor=\color{gray},
  frame=leftline,
  framexleftmargin=5pt
]{Chapters/04-FormalUsingAlloy/src/main.als}

  % 2 - Overall description
  \chapter{Overall description}
  \definecolor{gray}{HTML}{F7F6ED}

\lstinputlisting[
  language=alloy,
  numbers=left,
  stepnumber=1,
  breaklines=true,
  backgroundcolor=\color{gray},
  frame=leftline,
  framexleftmargin=5pt
]{Chapters/04-FormalUsingAlloy/src/main.als}

  % 3 - Specific requirements
  \chapter{Specific requirements}
  \definecolor{gray}{HTML}{F7F6ED}

\lstinputlisting[
  language=alloy,
  numbers=left,
  stepnumber=1,
  breaklines=true,
  backgroundcolor=\color{gray},
  frame=leftline,
  framexleftmargin=5pt
]{Chapters/04-FormalUsingAlloy/src/main.als}
  
  % 4 - Formal using Alloy
  \chapter{Formal using Alloy}
  \definecolor{gray}{HTML}{F7F6ED}

\lstinputlisting[
  language=alloy,
  numbers=left,
  stepnumber=1,
  breaklines=true,
  backgroundcolor=\color{gray},
  frame=leftline,
  framexleftmargin=5pt
]{Chapters/04-FormalUsingAlloy/src/main.als}

  % 5 - Effort spent
  \chapter{Effort spent}
  \definecolor{gray}{HTML}{F7F6ED}

\lstinputlisting[
  language=alloy,
  numbers=left,
  stepnumber=1,
  breaklines=true,
  backgroundcolor=\color{gray},
  frame=leftline,
  framexleftmargin=5pt
]{Chapters/04-FormalUsingAlloy/src/main.als}

  % 6 - References
  \begin{thebibliography}{0}
    \bibitem{latex}\href{https://www.latex-project.org/}{\LaTeX}  for writing the document {\scriptsize \newline$[$\url{https://www.latex-project.org/}$]$}

    \bibitem{alloy-style} \href{https://github.com/Angtrim/alloy-latex-highlighting}{Alloy Latex Highlighter}, a \LaTeX \ package to highlight alloy syntax that we modified to render improve the visualization of some operators
    {\scriptsize \newline$[$\url{https://github.com/Angtrim/alloy-latex-highlighting}$]$}

    \bibitem{title-pic} \href{http://tug.ctan.org/tex-archive/macros/latex/contrib/titlepic/titlepic.sty}{Title Pic}, a \LaTeX \ package that allows customization of the cover page
    {\scriptsize \newline$[$\url{http://tug.ctan.org/tex-archive/macros/latex/contrib/titlepic/titlepic.sty}$]$}

    \bibitem{draw-io} \href{https://www.draw.io/}{draw.io} to draw UML diagrams and sequence diagrams
    {\scriptsize \newline$[$\url{https://www.draw.io/}$]$}

    \bibitem{alloy-analyzer} \href{https://alloytools.org/}{Alloy Analyzer} to generate Alloy worlds and run alloy predicates
    {\scriptsize \newline$[$\url{https://alloytools.org/}$]$}

    \bibitem{tocbibind} \href{https://ctan.org/pkg/tocbibind}{tocbibind}, a \LaTeX \ package to bind the bibliography to the table of contents
    {\scriptsize \newline$[$\url{https://ctan.org/pkg/tocbibind}$]$}

    \bibitem{github} \href{https://github.com}{GitHub} for version control
    {\scriptsize \newline$[$\url{https://github.com}$]$}

    \bibitem{vscode} \href{https://code.visualstudio.com/}{Visual Studio Code} the IDE we used
    {\scriptsize \newline$[$\url{https://code.visualstudio.com/}$]$}

    \bibitem{wikibooks} \href{https://en.wikibooks.org/wiki/LaTeX/}{Wikibooks} we read the section about \LaTeX \ to learn how to structure the document
    {\scriptsize \newline$[$\url{https://en.wikibooks.org/wiki/LaTeX/}$]$}

    \bibitem{stackoverflow} \href{https://tex.stackexchange.com/}{Stack Exchange} for further doubts about \LaTeX
    {\scriptsize \newline$[$\url{https://tex.stackexchange.com/}$]$}

    \bibitem{Ionic} \href{https://ionicframework.com/}{Ionic} to generate mockups
    {\scriptsize \newline$[$\url{https://ionicframework.com/}$]$}

    \bibitem{GIMP} \href{https://www.gimp.org}{GIMP} to edit mockups
    {\scriptsize \newline$[$\url{https://www.gimp.org}$]$}
  \end{thebibliography}
\end{document}